\begin{afterword}
\summaryhead

The author often draws a scale of intelligence that runs from a ``village idiot'' to Einstein. That is the scale of the minds we meet. On the scale of possible minds the two are almost the same point, since both have the same brain design; a chimpanzee is much farther away, and the right side belongs to Jupiter Brains, Adams's Deep Thought, the Arisians and other minds of fiction. When Douglas Hofstadter said Einstein should be far to the right, the author saw a cultural gap. Reading science fiction, the author suspects, helps people imagine outside their own era, though this generalizes from fiction.

Human brains, adapted for hunting and for arguing over shares, use their evidence and their neurons inefficiently, and Einstein's feat was to repurpose one for physics. The author's childhood ideal was never Einstein but a dream of an AI. Ideals, Carl Schurz said, are like stars. An imagined ideal is only one's own mind talking and can lead astray, yet ideals made only of real, dead humans limit you to what has been done. Einstein talked nonsense about an impersonal God. The most important role models are dreams of perfection.

\responsehead

The post makes two claims: that the range of human minds is a small part of the range of possible minds, and that one should take one's ideals from imagined perfection rather than from great humans. The first is plausible and argued loosely. The second rests mostly on an image.

The scale is argued by brain design. The evidence given, that Einstein and the village idiot share a prefrontal cortex, a hippocampus and a cerebellum, does not separate them from a chimpanzee, which has the same parts. Brain size, which the post does not cite, does fit its ordering. Hofstadter's objection is reported in one sentence, and the ``cultural gap'' is the author's reading of it. The minds that fill the right of the scale come from fiction, which the post concedes. Its suspicion that science fiction widens the imagination sits uneasily with the author's earlier claim that other authors' imaginations stop people from using their own.

The question whether one superintelligence could do what the civilization of ``That Alien Message'' did gets an indirect answer. Biases, slow mental arithmetic and a fictional teacher's count of 5,760 insights a month show that human brains are inefficient. They do not show how large the gap is.

The case against human role models is that ``following stars'' at best gets you to the star. No case of someone held back this way is given, and Einstein's own record sits uneasily with it: he revered Newton and went past him. The one charge against Einstein's thinking outside physics, ``nonsense about an impersonal God,'' quotes nothing he said; his statements used ``God'' for the order of nature and rejected a personal God. The post admits its ideal ``harmed me to some degree'' and does not say how.

The post deserves credit for stating the objection to its own advice, that an imagined ideal is ``only your own mind talking,'' and for qualifying its claim about science fiction both ways. The figure of a hundred thousand years for our species was low even in 2008, but the point it serves survives.

In short: a plausible claim about the range of possible minds, supported by design talk and admitted fiction, and advice to aim at imagined perfection that rests on an image of stars and an unquoted dismissal of Einstein's views on God.
\end{afterword}
